\pagestyle{corporativo}

\begin{frontmatter}

	\title{\projecttitle}

	\author[inst1]{J. Arboleda\corref{cor1}}
	\ead{[correo@dmlsas.com]}

	\author[inst1]{F. Navia}

	\author[inst1]{H. Rosero}

	\cortext[cor1]{Autor correspondiente}
	\address[inst1]{DML Ingenieros Consultores S.A.S., Cali, Colombia, 760025}

	\begin{abstract}
		\thispagestyle{corporativo}
		Se presenta el an\'{a}lisis de flexibilidad de la tuber\'{i}a de descarga
		de la bomba de recirculaci\'{o}n 2 del sistema DDW (l\'{i}nea SIM-003,
		job CAESAR P2603-PR-SIM-003) del proyecto P2603 SW-K60 para Smurfit
		Westrock, ejecutado conforme a ASME B31.3-2016 con el software
		CAESAR II 2019. La l\'{i}nea, construida en acero inoxidable
		ASTM A312 TP 304L Sch 10S en di\'{a}metros de 1.5 y 6~NPS, opera con
		agua a 1241~\si{kPa}~g y 90~\si{\celsius}. El modelo se gener\'{o}
		desde el isom\'{e}trico de AutoCAD Plant 3D 2027 mediante archivo PCF
		corregido, y se evaluaron nueve casos de carga (operaci\'{o}n,
		sostenidos y expansi\'{o}n). El esfuerzo m\'{a}ximo de c\'{o}digo fue
		\num{18089.2}~\si{kPa} frente a un admisible de
		\num{115142.4}~\si{kPa}, correspondiente al nodo 120 en el caso
		sostenido alterno 4 (Alt-SUS) W+P2, con una relaci\'{o}n de
		utilizaci\'{o}n de 15.7~\%. La evaluaci\'{o}n de cumplimiento de
		c\'{o}digo arroj\'{o} \textbf{CODE COMPLIANCE EVALUATION PASSED}: la
		l\'{i}nea es aceptable seg\'{u}n ASME B31.3-2016 con amplio margen
		respecto al esfuerzo admisible.
	\end{abstract}

	\begin{keyword}
		flexibilidad de tuber\'{i}as \sep CAESAR II \sep ASME B31.3 \sep
		esfuerzo sostenido \sep desplazamientos \sep an\'{a}lisis de esfuerzos
	\end{keyword}

\end{frontmatter}

\pagestyle{corporativo}
\thispagestyle{corporativo}
